\paragraph{Counterfactual-environment prompting.} We next tested whether explicitly asking GPT-5.5 to reason about
hypothetical changes to the store (rearranged or rotated shelves, moved stations and exits, changed routes) helps it
find layout-independent norms when it has seen only one store. Apart from this guidance paragraph, the prompts were
identical to the previous experiment. With single-layout demonstrations, the intended norm appeared in GPT-5.5's
candidates in 6/10 runs, the same rate as with the earlier generalization guidance (6/10, on different runs).
Coordinate-based proposals were absent under both. Selection did not improve: in every run the Bayesian learner chose
a shorter coordinate rule consistent with the demonstrations, and accuracy on an unseen test layout remained at chance
(0.50, violation F1 0.00). With demonstrations from two layouts, results matched the earlier guidance (intended norm
selected in 5/10 runs; test-layout accuracy 0.90, F1 0.92). The remaining errors again came from over-broad relational
definitions, despite the prompt's explicit warning against them. Counterfactual reasoning in the prompt therefore did
not substitute for observing a second environment. The bottleneck was not proposal but the evidence available to
choose between hypotheses that explain the observed store equally well. Because the two-layout condition also has
twice as many demonstrations, these results do not separate layout diversity from demonstration count.
